\chapter{Introducción}

Un acelerador lineal de partículas (linac) es un dispositivo
que permite incrementar la energía cinética de partículas
cargadas mediante campos eléctricos de radiofrecuencia (RF),
dispuestos a lo largo de una estructura de aceleración.
En los aceleradores lineales de tubos de deriva, particularmente
en la estructura tipo Alvarez, las partículas atraviesan
sucesivamente regiones de aceleración denominadas huecos,
en las cuales el campo eléctrico longitudinal les transfiere
energía. Entre estos huecos se encuentran los tubos de deriva,
que protegen a las partículas del campo eléctrico oscilante
durante su tránsito. De esta manera, la variación temporal
de la fase del campo permite que las partículas alcancen
los siguientes huecos en una fase adecuada para continuar
su aceleración \parencite{wangler,wille}.

El funcionamiento de este tipo de aceleradores requiere
mantener el sincronismo entre el movimiento de las partículas
y la oscilación del campo de radiofrecuencia. Para ello,
la longitud de las celdas debe ajustarse a la velocidad
de la partícula síncrona. A medida que esta incrementa
su energía cinética, también aumenta su velocidad, por
lo que las celdas deben adquirir longitudes progresivamente
mayores para conservar la condición de sincronismo.
Este principio constituye una característica fundamental
de los aceleradores lineales de tubos de deriva.

Los aceleradores lineales tienen diversas aplicaciones
en física y tecnología. Se emplean como etapas de inyección
en instalaciones de aceleración de partículas, en equipos
de radioterapia que utilizan haces de electrones con energías
típicamente comprendidas entre
$6\,\mathrm{MeV}$ y $18\,\mathrm{MeV}$,
y en instalaciones que producen neutrones mediante procesos
de espalación. Estas aplicaciones evidencian la importancia
de comprender los mecanismos de transferencia de energía
y las condiciones dinámicas que permiten mantener
la estabilidad del haz durante el proceso de aceleración.

En este contexto, el presente trabajo desarrolla un modelo
computacional de un acelerador lineal de protones con
estructura tipo Alvarez, operando a una frecuencia de
$352.2\,\mathrm{MHz}$, para estudiar la aceleración desde
$3\,\mathrm{MeV}$ hasta $50\,\mathrm{MeV}$.
El estudio comprende el diseño de las longitudes de las
celdas para la partícula síncrona, el análisis de la dinámica
longitudinal de un paquete de partículas y la evaluación
de la estabilidad de fase. Asimismo, se examinan las
oscilaciones de sincrotrón, el efecto de las desviaciones
en la frecuencia de radiofrecuencia y el comportamiento
de los electrones en el régimen ultrarrelativista.
Mediante la implementación numérica en Python, se busca
analizar la evolución del sistema y contrastar los resultados
de la simulación con las predicciones del modelo teórico.
